% introduction  -  part of main.tex (see journal/README.md)
A transformer carries no state from one call to the next. Everything it knows about a
conversation has to be supplied again in its prompt, the window of text it reads at once
\cite{vaswani2017attention}. Systems that give language models long-term memory therefore
treat memory as a layer of engineering wrapped around a frozen network: text is embedded,
stored in an external database, retrieved by similarity and pasted back into the context
window\cite{lewis2020rag,packer2023memgpt,park2023generative}. That works, but it ties
recall to the context budget, leaves the model itself stateless, and hands memory to the
harness instead of the model. Memory-augmented architectures move memory inside the network
instead---as recurrent or compressed
state\cite{dai2019transformerxl,rae2020compressive,bulatov2022rmt}, as a memory bank read
and written by attention\cite{graves2014ntm,graves2016dnc,weston2015memnet,sukhbaatar2015e2e,wu2022memformer},
or as a cache of past attention keys searched by nearest neighbour\cite{wu2022memorizing}.
That memory is bound to one sequence, is addressed by contextual representations that drift
during training, and is not a portable artefact of the model.

We ask whether memory can be part of the model. We add one memory block inside the stack
(Fig.~\ref{fig:arch}a). Every layer reads from, and writes its own state into, a single
persistent store shared across layers. The store's address is the raw embedding of each
token (a word or word fragment) at layer~0, before any processing and without position
information, so retrieval is an exact token match. Text is processed in fixed segments,
and the model keeps no key--value cache (the attention state a
transformer would otherwise carry forward) between them, so the store is the only channel
from one segment to the next.

The design puts the weights and the store in separate roles. The weights are stateless and
the store is the state. A frozen model carries nothing of a conversation; all of it lives in
a detachable file. We show this literally: swapping the store swaps the memory without
changing a weight (Fig.~\ref{fig:arch}b). We also show where the property stops. A store can
be read only by the model that wrote it.

Across 113 controlled experiments, we find that this block
(i)~recalls facts stated once in natural-language text and queried entirely out of context,
and keeps them across a process restart; (ii)~holds facts far beyond the attention window,
read by touching a single slot, with recall flat to a million tokens and no lost-in-the-middle
effect (the weaker retrieval of facts placed mid-context\cite{liu2024lost}); (iii)~can be
revised, bounded and reasoned over, which turns a lookup table into a memory; and
(iv)~maps onto the shared key--value constructions now used to compress attention
caches\cite{deepseekv41flash2026}. We also locate where it fails. A single store saturates
near 128 distinct facts, and a fact written late into a saturated store is lost. Both limits
come from addressing and from the write policy. Position is not the problem.